\section{Empirical implications}\label{sec:empirical}

The model gives predictions for publicly visible sale processes in which one bidder is prepared and a further buyer
has not yet committed to diligence. This section states them for the benchmark auction at the prior \(1/2\), with a
slot in the band of Section~\ref{sec:benchmark}. Appendix~\ref{app:empirical} states each result in full and proves
it. Full two-sided trade is the unique outcome at every strength \(r\) with \(r_S<r<h\), where \(r_S\) solves
\(\rho\DT(r)m=k\) exactly. This interval contains \(r_1\), \(r_C\) and \(r_2\). Under full two-sided trade, the
informed share of nonzero orders is \(\vartheta=u/(2+u)\).

\paragraph{Entry and incumbent strength.} Proposition~\ref{prop:ccc} and Proposition~\ref{prop:control} compare
processes with prepared bidders of different public strength. When the price reaches the further buyer before its
decision, a stronger prepared bidder can raise the probability that the buyer prepares. In the benchmark, it rises by
\(\reg{entry_jump_B}\) from \(r_0\) to \(r_1\). When the price does not reach the buyer, a stronger prepared bidder
does not raise it. Against a strong enough prepared bidder, preparation returns to its floor while the price stays
informative.

\paragraph{The spread and its entry component.} The spread is the difference between the conditional execution
prices of a unit purchase and a unit sale. It has a class component \(\rho\DT\vartheta\) and an entry component, which
is positive exactly when a purchase recruits (Proposition~\ref{prop:spread}). The sale impact contains only half of the
class component. The purchase impact therefore exceeds the sale impact exactly when a purchase recruits the buyer, and
this test needs no knowledge of the auction objects.

In the benchmark at \(r_1\), the spread is \(\reg{ba_r1_B}\), and
its entry component is \(\reg{ba_share_r1_B}\) of it (computer-assisted). With a hidden price, the spread would be
\(\reg{ba_class_r1_B}\). At \(r_2\) the spread is \(\reg{ba_r2_B}\), the class component alone, and at \(r_C\) it falls
by \(\reg{ba_jump_rC_B}\).

\paragraph{The informed share.} Order frequencies alone leave the slot open. Under full two-sided trade, the order law
depends on \((\lambda,\beta)\) only through the volume of nonzero orders (Proposition~\ref{prop:informedshare}). The
investor's incentives limit the shares that the counts allow. At \(r_1\), they allow every informed share from
\(\approx\reg{theta_min_counts_B}\) up to, but not including, \(\reg{theta_max_full_eq_r1_B}\).

A slot with share \(\reg{alt_theta_B}\) has the
benchmark volume and the benchmark entry, but its spread at \(r_1\) is \(\reg{alt_ba_r1_B}\), against \(\reg{ba_r1_B}\).
Prices therefore separate slots that counts cannot. Given the entry floor and the target-payoff spread, the sale impact
and the volume identify the informed share.

\paragraph{Winners.} Among prepared challengers, the share of the high class at each price is the price-implied
posterior, which rises with the price (Proposition~\ref{prop:winner}). A prepared high-class challenger always outbids
the incumbent, and a low-class one loses whenever the incumbent's value exceeds \(\ell\). Orders therefore predict the
winner even where they no longer predict entry. In the benchmark, a prepared challenger wins more often
after a purchase than after a sale, by \(\reg{wgap_r1_B}\) at \(r_1\) and by \(\reg{wgap_r2_B}\) at \(r_2\). The
covariance of the order with entry is \(\reg{cov_entry_r1_B}\) at \(r_1\) and zero at \(r_2\), while its covariance with
a challenger win stays positive.

\paragraph{Announcements.} Add a public event after the price that reveals only whether the challenger prepared. The
mean price change over both reports, with their model probabilities, is zero. The mean change in a sample of
preparation reports alone is positive (Proposition~\ref{prop:announcement}): \(\reg{evprep_mean_r1_B}\) at \(r_1\) in
the benchmark. Below the ceiling, an order opportunity with no order is bad news about entry and no news about the
class. At \(r_1\), the price after no order minus the expected price before the opportunity is
\(\reg{pdrop_neutral_r1_B}\).

\paragraph{Two clusters at intermediate strengths.} On \((r_F,r_N)\), Section~\ref{sec:coexistence} gives two
equilibria under rule C1. In one, there is no informed trade, the mean order is zero and the three prices coincide. In
the other, the high class buys and the low class abstains. The mean order is then \(\lambda/2\)
(\(\reg{imb_buy_B}\) in the benchmark), and every challenger prepares after a purchase. A sale executes above the price
after no order, by \(\reg{U_gap_r160_B}\) at \(\reg{str_r160}\).

Under rule C1, processes at one public strength in this
range therefore form two clusters. Under rule CM, the threshold equilibrium of Remark~\ref{rem:ties} adds a third
pattern. Which equilibrium a process selects is open.

\paragraph{Measurement.} The counterpart of preparation is the start of substantive diligence or a proposal that
requires costly preparation; public offers are a narrower outcome. Incumbent strength is measured from information
that is public before that decision, and the price must precede it.

The order is the first prespecified public order opportunity after the sale becomes public, and an opportunity with
no trade is recorded as a zero. A sample of first
executed trades omits the zero order, so it cannot identify the volume, the effect of an inactive opportunity, or the
informed share. The announcement formulas need events that reveal only preparation. Target returns before a bid also
partly reflect anticipation of that bid \citep{Schwert1996}, so each test dates the price before the buyer's decision. The trading counterpart is outside trade with a view
on the further buyer's fit, such as flow from sector specialists.
